\textbf{Pour cette première partie, aucune justification n'est demandée et une seule réponse est possible par question.} Les justifications ci-dessous sont fournies à titre pédagogique.

Une réponse fausse ou l'absence de réponse n'enlève aucun point : il ne faut donc jamais laisser une question sans réponse.

\subsection*{Question 1 \quad Réponse D}

Une réduction de 25\,\% correspond au coefficient multiplicateur
$1 - \dfrac{25}{100} = 0,75$ :
\[80 \times 0,75 = 60.\]
Le nouveau prix du pantalon est donc de 60~\euro.

\subsection*{Question 2 \quad Réponse C}

Le loyer représente 30\,\% du salaire s : $600 = 0,30 \times \text{s}$, donc
\[ \text{s} = \dfrac{600}{0,30} = \np{2000}.\]
Le salaire de Paul est de \np{2000}~\euro.

\subsection*{Question 3 \quad Réponse B}

La fréquentation de juillet (532 spectateurs) correspond à celle de juin (notée juin), diminuée de 8\,\%. Le coefficient multiplicateur associé est $1-0,08=0,92$.

On a
$532 = \text{juin} \times 0,92$, d'où
\[\text{juin} = 532 \div 0,92.\]

\subsection*{Question 4 \quad Réponse C}

On réduit les quatre fractions au même dénominateur (120) pour les comparer :
\[\dfrac{1}{4} = \dfrac{30}{120}, \quad \dfrac{3}{4} = \dfrac{90}{120}, \quad
\dfrac{1}{8} = \dfrac{15}{120}, \quad \dfrac{2}{15} = \dfrac{16}{120}.\]
La plus petite est donc $\dfrac{1}{8}$.

\subsection*{Question 5 \quad Réponse A}

On calcule chaque écriture :
\[\dfrac{\np{1025}}{100} = 10,25, \qquad \dfrac{205}{20} = 10,25, \qquad 102,5 \times 10^{-1} = 10,25,\]
mais
\[\np{0,1025} \times 10^{-2} = \np{0,001025} \neq 10,25.\]
C'est donc la première écriture qui est différente des trois autres.

\subsection*{Question 6 \quad Réponse A}

Sur le graphique, chaque activité est associée à deux pourcentages : celui des sportifs de
moins de 40 ans et celui des sportifs de 40 ans et plus. Parmi les valeurs de la série \og moins de 40 ans \fg, seule celle du VTT vaut
11\,\%.

\subsection*{Question 7 \quad Réponse B}

Le tableau de signes indique que $f$ s'annule en 3, est négative avant et
positive après : le coefficient directeur de $f$ est donc positif.
Seule la proposition $f(x) = 2x - 6$ vérifie ces deux conditions : elle s'annule bien en 3
(puisque $2 \times 3 - 6 = 0$), et son coefficient directeur 2 est positif.

\subsection*{Question 8 \quad Réponse D}

La courbe monte de $x=-4$ (où $f(-4) = -5,4$) jusqu'à un maximum local en $x=-1$ (où
$f(-1) = 3,6$) : $f$ est donc croissante sur $[-4 ; -1]$.

\subsection*{Question 9 \quad Réponse D}

On isole $c^2$ dans la formule, puis on prend la racine carrée (une vitesse est positive) :
\[E = mc^2 \iff c^2 = \dfrac{E}{m} \iff c = \sqrt{\dfrac{E}{m}}.\]

\subsection*{Question 10 \quad Réponse C}

$289 \approx 300$ et $0,307 \approx 0,3$, donc
\[289 \times 0,307 \approx 300 \times 0,3 = 90.\]
Seule la réponse $88,723$ est de cet ordre de grandeur.

\subsection*{Question 11 \quad Réponse A}

On applique l'identité remarquable $(a-b)^2 = a^2 - 2ab + b^2$ avec $a = 2x$ et $b = 7$ :
\[(2x-7)^2 = (2x)^2 - 2 \times 2x \times 7 + 7^2 = 4x^2 - 28x + 49.\]

\subsection*{Question 12 \quad Réponse B}

On convertit les heures en minutes en multipliant par 60 :
\[1,2 \times 60 = 72.\]
1,2~h correspond donc à 72~minutes.
